\documentclass[a4paper,12pt,twoside,openright,titlepage]{book}

%Additional packages
\usepackage[utf8]{inputenc}
\usepackage[T1]{fontenc}
\usepackage[dutch,english]{babel}
\usepackage{imakeidx}
\usepackage{syntonly}
\usepackage[official]{eurosym}
%\usepackage[graphicx]
\usepackage{graphicx}
\graphicspath{ {./images/} }
\usepackage{float}
\usepackage{xurl}
\usepackage{hyperref}
\hypersetup{colorlinks=true, linkcolor=blue, citecolor=blue, filecolor=blue, urlcolor=blue, pdftitle=, pdfauthor=, pdfsubject=, pdfkeywords=}
\usepackage{tabularx}
\usepackage[table]{xcolor} % Table colors
\usepackage{scrextend}
\addtokomafont{labelinglabel}{\sffamily}
\usepackage{listings}
\usepackage{adjustbox}
\usepackage{color}
\usepackage{csquotes}

% Define colors
\definecolor{ashgrey}{rgb}{0.7, 0.75, 0.71}

% Listing style
\lstset{
  backgroundcolor=\color{ashgrey}, % choose the background color; you must add \usepackage{color} or \usepackage{xcolor}; should come as last argument
  basicstyle=\footnotesize,        % the size of the fonts that are used for the code
  breakatwhitespace=true,          % sets if automatic breaks should only happen at whitespace
  breaklines=true,                 % sets automatic line breaking
  extendedchars=true,              % lets you use non-ASCII characters; for 8-bits encodings only, does not work with UTF-8
  frame=single,	                   % adds a frame around the code
  rulecolor=\color{black},         % if not set, the frame-color may be changed on line-breaks within not-black text (e.g. comments (green here))
  keepspaces=true,                 % keeps spaces in text, useful for keeping indentation of code (possibly needs columns=flexible)
  columns=fullflexible,		   % make copy and paste possible
  showstringspaces=false,          % if true show spaces in strings adding particular underscores
  showspaces=false,                % if true show spaces everywhere adding particular underscores; it does not override 'showstringspaces'
}

% Uncomment for production
% \syntaxonly

% Style
\pagestyle{headings}

% Turn on indexing
\makeindex[intoc]

% Define document
\author{D. Leeuw}
\title{LDAP}
\date{\today\\v.0.4.0}

\begin{document}
\selectlanguage{dutch}

\maketitle

\copyright\ 2023 Dennis Leeuw\\

\begin{figure}
\includegraphics[width=0.3\textwidth]{CC-BY-SA-NC.png}
\end{figure}

\bigskip

Dit werk is uitgegeven onder de Creative Commons BY-NC-SA Licentie en laat anderen toe het werk te kopi\"eren, distribueren, vertonen, op te voeren, en om afgeleid materiaal te maken, zolang de auteurs en uitgever worden vermeld als maker van het werk, het werk niet commercieel gebruikt wordt en afgeleide werken onder identieke voorwaarden worden verspreid.


%%%%%%%%%%%%%%%%%%%
%%% Introductie %%%
%%%%%%%%%%%%%%%%%%%

\frontmatter
\chapter{Over dit Document}
\section{Leerdoelen}
Dit document gaat over ethernet. Het leert je wat ethernet is. Om meer te weten te komen over wat er op je netwerk gebeurt maak je ook kennis met de tools: \texttt{arp} en \texttt{wireshark}.

Na het bestuderen van dit document heb je een basis kennis van ethernet, weet je wat het adres-formaat is van ethernet en kan je informatie over het netwerk opvragen via verschillende tools.

In de praktijk opdrachten leer je om je eigen kabel te maken en deze te testen. Tevens maak je je eigen Wi-Fi netwerk.

\section{Voorkennis}
Er wordt van de lezer verwacht dat deze kennis heeft van de volgende onderwerpen:
\begin{itemize}
\item TCP/IP
\item CMOS en BIOS tijd
\end{itemize}


%%%%%%%%%%%%%%%%%
%%% De inhoud %%%
%%%%%%%%%%%%%%%%%
\tableofcontents

\mainmatter
\chapter{Inleiding}
Bestandssystemen gebruiken nummers om bestanden te koppelen aan gebruikers. Bij een bestand wordt de User-ID en de Group-ID opgeslagen om te bepalen welke gebruikers en groepen toegang hebben tot een bestand. Als we een bestandssysteem willen delen over het netwerk om data met elkaar te kunnen delen dan moeten alle systemen die bij de gedeelde data kunnen weten welke UID of GID bij welke gebruikers en groepen horen. We kunnen natuurlijk proberen om dat op elk systeem handmatig op orde te maken, maar het zou makkelijker zijn als we een centrale database hebben waar elke machine aan kan vragen wat er bij een UID of GID hoort.

Een mogelijke centrale netwerk oplossing die heel veel gebruikt wordt is LDAP\index{LDAP}. De afkorting LDAP staat voor Lightweight Directory Access Protocol\index{Lightweight Directory Access Protocol}. Het Lightweight deel staat voor het feit dat het een uitgeklede versie is van het ITU X.500 protocol. X.500, afkomstig uit de telefoniewereld, kent het DAP (Directory Access Protocol) om data uit het systeem te halen. DAP is heel uitgebreid en dus arbeidsintensief, vandaar dat er voor gebruik over TCP/IP netwerken een lichter alternatief werd bedacht: LDAP.

LDAP is een netwerk protocol dat beschrijft hoe data aangeboden moet worden aan het LDAP systeem. De data die uit het systeem verkregen worden of erin gestopt worden zijn in het LDIF formaat.



\chapter{Directory Information Tree}
LDAP\index{LDAP} is een database met informatie over een gebruiker, zoals de gebruikersnaam, de groepen waarvan de gebruiker lid is en de User ID, maar ook gegevens als telefoonnummers, e-mail adressen en andere gegevens van een gebruiker kunnen worden opgenomen in LDAP. We kunnen LDAP dan ook vergelijken met een adresboek (telefoonboek). In LDAP kan ook data worden opgeslagen over computers en groepen. Kortom LDAP is eigenlijk een manier om data gestructureerd op te slaan. De structuur is vergelijkbaar met een bestandssysteem, een hi\"erarchische manier, met mappen en submappen waarin data staat. Vandaar de naam Directory\index{Directory}\index{LDAP!Directory} in LDAP.

Een \textquote{map} in LDAP zou een gebruiker kunnen zijn waaronder de gegevens van de gebruiker worden opgeslagen. Een map kan ook een bedrijf zijn met daarin de bedrijfsgegevens, zoals adressen, met daaronder mappen voor werknemers met daarin de gebruikersgegevens en contactgevens van deze medewerkers. Een LDAP-tree kan dus op verschillende manieren gestructureerd zijn. Via zoekopdrachten kan er in deze databoom gezocht worden. Binnen LDAP heten de mappen objecten\index{Object}. Een object kan andere objecten bevatten.

LDAP is een hi\"erarchische tree van opgeslagen data. Via zogenaamde schema bestanden wordt bepaald wat er in een object in de boom opgeslagen kan, mag of moet worden. Een schema bestand zegt dus wat bij bijvoorbeeld een user object er aan data aanwezig moet zijn en welke data er ook aanwezig mag zijn. De gebruikersnaam is verplicht, net als een UID, een telefoonnummer is niet verplicht. Dit geldt voor een LDAP systeem dat is ingericht met schema bestanden voor gebruik als gebruikersdatabase.

Omdat elke organisatie zijn eigen LDAP-tree kan hebben met zijn eigen gebruikers (zijn eigen interne adresboek) en LDAP systemen theoretisch wereldwijd te koppelen zijn is het handig om de basis van de tree (de root) aan te laten sluiten bij wereldwijd unieke naam van de organisatie. De keuze binnen de meeste organisaties is gevallen op het aansluiten bij het DNS systeem, daar dit het enige systeem is dat een unieke naam heeft voor een organisatie namelijk zijn domeinnaam. Er kan wereldwijd maar \'e\'en microsoft.com zijn.

LDAP werk niet met domeinnamen, maar met objecten. Waarbij het eerst genoemde object de root. Zo zou de basis van de LDAP tree voor Microsoft DC=microsoft,\allowbreak DC=com kunnen zijn. DC\index{DC} staat voor Domain Component\index{Domain Component} en verwijst naar het feit dat deze LDAP tree is opgebouwd volgens het DNS systeem. In de notatie valt op dat het aansluit op de DNS structuur met een beschrijving van elk object in de tree. Zo zouden we onze gebruikers kunnen stoppen in een organisatorische eenheid ofwel een Organizational Unit\index{Organizational Unit}. Laten we daarvoor de LDAP eenheid nemen: \index{OU}OU=users,\allowbreak DC=microsoft,\allowbreak DC=com. Daaronder kunnen we onze gebruikers stoppen CN=bgates,\allowbreak OU=users,\allowbreak DC=microsoft,\allowbreak DC=com. CN\index{CN} staat hierbij voor Common Name\index{Common Name}. Sommige organisaties gebruiken voor de Common Name de loginnaam, andere gebruiken de volledige naam van de gebruiker of het nummer van werknemer. Het is belangrijk dat de naam van elk object binnen een tak (branch) van de boom uniek is. Er kan dus niet nog een CN=bgates voorkomen binnen OU=users,\allowbreak DC=microsoft,\allowbreak DC=com. Een object in de boom die andere objecten kan bevatten worden ook wel Containers\index{Container} genoemd, in de Windows wereld wordt met Container vaak een Common Name object bedoeld.

Het \textquote{pad} CN=bgates,\allowbreak OU=users,\allowbreak DC=microsoft,\allowbreak DC=com is een unieke entry van een object in LDAP. Dit noemen we de Distinguished Name\index{Distinguished Name}. De Distinguished Name is het volledige pad naar een object in de tree.

De volledige LDAP data boom van een organisatie noemen we een DIT\index{DIT} (Directory Information Tree\index{Directory Information Tree}). Een DIT bevat zoals gezegt objecten\index{Object} en objecten hebben Attributes\index{Attribute}. Een attribute kan bijvoorbeeld het telefoonnummer zijn.



\chapter{LDIF}

%%%%%%%%%%%%%%%%%%%%%
%%% Index and End %%%
%%%%%%%%%%%%%%%%%%%%%
\backmatter
\printindex
\end{document}

%%% Last line %%%
